\documentclass[book.tex]{subfiles}

\begin{document}

%\pagestyle{empty}
%\begin{center}\Huge\bfseries\itshape
\part{Foundational Physics}
\label{part:physics}
%\end{center}

In \autoref{part:math} of this book, we reviewed the mathematical foundations necessary to understand physics from a differentiable programming viewpoint. 

This part will introduce you to the foundations of physics by studying a variety of different physical systems using the mathematical machinery we have built up in \autoref{part:math} alongside the type-theoretic machiner from \autoref{part:diffprog}. Our main goal in this part of the book will be to explain how to use the tools of type theory and programming to model physical systems.

If you are a physicist, you might find a lot of the physics here basic, but you might perhaps find our exposition novel and worth a skim. If you're a computer scientist or applied scientists, our exposition will hopefully help you get a different way to understand the foundations of physics by using types and programming.

In the next part of the book, we will combine our type-driven physical machinery with differentiable programming to demonstrate how differentiable physics can provide deeper understanding of the physical systems we study in the Part.




\end{document}